% ============================================================
%  INTRODUCTION
% ============================================================


Spent nuclear fuel discharged from conventional light water reactors contains significant quantities of plutonium generated in-core as a byproduct of neutron capture in $^{238}$U. During reactor operation, the abundant $^{238}$U in low-enriched uranium (LEU) fuel undergoes the transmutation chain $^{238}$U$(n,\gamma) \to$ $^{239}$U $\to$ $^{239}$Np $\to$ $^{239}$Pu, accumulating reactor-grade plutonium throughout irradiation. Even in spent fuel form where mixed with fission products and higher actinides --- this plutonium represents a long-term proliferation concern: under treaty withdrawal scenarios, a state with access to civilian spent fuel pools could in principle separate weapon-usable material via established chemical processes 
\cite{Feiveson2004NuclearDiversion}.

The predominant strategies for managing this stockpile carry their own liabilities. Chemical reprocessing programs (PUREX) employed in France, Russia, Japan, and the United Kingdom recover fissile material for Mixed Oxide (MOX) fuel fabrication, but the process necessarily generates streams of separated plutonium, elevating diversion risk despite extensive international safeguards \cite{2008SpentIAEA-TECDOC-1587}. The Direct Use of Spent PWR Fuel In CANDU (DUPIC) process avoids chemical separation by mechanically treating and re-pelletizing spent fuel, but requires decladding, oxidation--reduction cycles, powder blending, and refabrication in heavily shielded hot cells --- a process that costs on the order of \$560/kgHM in refabrication alone and demands substantial specialized infrastructure \cite{Ko2001EconomicCost}.

Beyond proliferation risk and reprocessing economics, the accumulated spent fuel stockpile carries direct implications for long-term energy security. Between the years 2003–2013, the United States alone accumulated $\sim$15,300 metric tonnes of spent PWR fuel, with an average discharge burnup of 47.08 GWd/tHM and initial enrichment of 4.28\% \cite{nureg_cr7227}. When loaded with similarly characterized fuel, \textit{Last Breath Reactor} (LBR) has demonstrated extending burnup to 61.1 GWd/tHM while remaining critical. In practical terms, this means that a fleet of LBRs could produce $\sim$2,559 GW-years (thermal) of energy from the spent fuel collected within that 10-year period.

This inventory represents both a regulatory burden and a latent energy resource: the residual fissile and fertile content of spent LWR assemblies remains substantial relative to the additional energy that could in principle be extracted under extended burnup conditions. An upward revision of licensed burnup limits on the order of $+10$\,GWd/t --- for example, from a typical limit of $\sim$62\,GWd/tU to $\sim$72\,GWd/tU --- would meaningfully reduce the number of assemblies discharged per unit of electrical energy generated across the operating fleet of approximately 90 commercial reactors. Compounded over decadal timescales, this translates into deferred pressure on interim storage capacity, reduced dry cask procurement, and lower cumulative transportation requirements. More broadly, increased burnup improves the energy density of the domestic fuel cycle, reducing dependence on uranium import supply chains and enhancing the strategic resilience of nuclear generation relative to fuel sources subject to geopolitical disruption.


The \textit{Last Breath Reactor} (LBR) is designed to address the plutonium stockpile problem through a fundamentally different approach: the direct, unmodified reuse of spent LWR fuel assemblies as reactor fuel, with no decladding, no chemical separation, and no pellet refabrication. Spent fuel rods discharged from conventional PWRs are inserted directly into the LBR core, where a thermalized heavy water neutron spectrum is engineered to fission $^{239}$Pu faster than it is bred from $^{238}$U. This approach eliminates the proliferation-sensitive processing steps inherent to both PUREX and DUPIC while simultaneously extracting additional commercial-scale electrical energy from material already regarded as waste.


The central technical challenge is maintaining criticality in a reactor loaded with highly-depleted spent fuel ($\sim$52\,GWd/tHM initial burnup) while sustaining a neutron spectrum sufficiently thermal to achieve net plutonium burning. Spent fuel from the final irradiation cycle contains a degraded fissile inventory, elevated fission product poisoning, and a spatially non-uniform burnup distribution along each rod --- the rod extremities retaining higher fissile concentrations than the axial midplane due to flux depression by structural hardware and control rod shadowing. The LBR design must accommodate these heterogeneities while meeting commercial power output, regulatory burnup limits, and thermal safety requirements.

This report documents the full design process from initial concept through high-fidelity analysis, including: a theoretical basis for the transmutation strategy, preliminary 0-D and infinite-lattice scoping calculations, Serpent~2 Monte Carlo depletion analyses benchmarked against measured spent fuel data, design iteration toward a novel power-redistribution device, and thermal-hydraulic and structural qualification of the selected configuration.

